% ============================================================
% appendix/prompt_floats.tex
% 浮动体部分：DS/RF system prompt 各自独立成图（含独立 caption）
% + 七个任务的完整 user prompt。
% 排版顺序：先 RU 的四个图（两两合并到一页），再 RW 的四个图
% （Coherence Part 1 / Part 2 / Positioning Type / Positioning Check，
%  其中 Positioning Check 放在最后）。
% 依赖 appendix/prompts.tex 中的定义（main.tex 已保证先加载该文件）。
% 所有 user prompt 中 [QUERY]: / [CRITERIA]: / [ANSWER]: 均已加粗。
% ============================================================

% ---------- Shared system prompts ----------
% The two system prompts are shown as two separate figures.  The RF
% wording below uses the no-comma form; the released RU records contain
% the otherwise identical ``Then, output'' spelling.
\clearpage
\begin{figure*}[t]
\begin{promptboxbig}{Direct-score (DS) - system prompt}
  \ptextsys{\SystemScoreOnly}
\end{promptboxbig}
\caption{The shared DS system prompt used by the score-only branch,
common to all seven in-domain tasks.}
\label{fig:so}
\end{figure*}

\begin{figure*}[t]
\begin{promptboxbig}{Rationale-first (RF) - system prompt}
  \ptextsys{\SystemCoTThen}
\end{promptboxbig}
\caption{The shared RF system prompt used by the rationale-first branch,
common to all seven in-domain tasks. Related Work records use ``Then
output'', while the four Review Utility records use the otherwise
identical ``Then, output'' wording.}
\label{fig:rs}
\end{figure*}

% ---------- Complete user prompts for the seven in-domain tasks ----------
% Each panel contains the complete user message from the first released test
% instance for that task, including its real ANSWER field.  System prompts and
% assistant completions are intentionally omitted.

% ---------- RW -- Coherence (split into two parts) ----------
\newcommand{\UserCoherenceRealPartOne}{%
\textbf{[QUERY]:} Given a context from a scientific paper, your task is to write a coherent citation sentence that cites the given paper with a specific citation number. Citation number will be also provided along with the context.


\textbf{[CRITERIA]:} The paper context supports (entails) the citation sentence. In cases where more than one paper is referenced in the sentence, as long as context in which given citation number fits the paper content, it should be count as entailment as well. In multiple citation cases, the paper does not have to entail whole sentence. Scoring rubric is as follows:
0: The given context does not align with the context in which the citation number appears in the sentence, or with the overall meaning of the sentence itself.
1: The given context is compatible with the context in which the citation number appears in the sentence.





[CONTEXT]: The thud of a bouncing ball, the onset of speech as lips open -- when visual and audio events occur together, it suggests that there might be a common, underlying event that produced both signals. In this paper, we argue that the visual and audio components of a video signal should be modeled jointly using a fused multisensory representation. We propose to learn such a representation in a self-supervised way, by training a neural network to predict whether video frames and audio are temporally aligned. We use this learned representation for three applications: (a) sound source localization, i.e. visualizing the source of sound in a video; (b) audio-visual action recognition; and (c) on/off-screen audio source separation, e.g. removing the off-screen translator's voice from a foreign official's speech. Code, models, and video results are available on our webpage: this http URL
As humans, we experience our world through a number of simultaneous sensory streams. When we bite into an apple, not only do we taste it, but -as Smith and Gasser [1] point out -we also hear it crunch, see its red skin, and feel the coolness of its core. The coincidence of sensations gives us strong evidence that they were generated by a common, underlying event [2] , since it is unlikely that they co-occurred across multiple modalities merely by chance. These cross-modal, temporal co-occurrences therefore provide a useful learning signal: a model that is trained to detect them ought to discover multimodal structures that are useful for other tasks. In much of traditional computer vision research, however, we have been avoiding the use of other, non-visual modalities, arguably making the perception problem harder, not easier.
}

\newcommand{\UserCoherenceRealPartTwo}{%
\textbf{[QUERY]:} Given a context from a scientific paper, your task is to write a coherent citation sentence that cites the given paper with a specific citation number. Citation number will be also provided along with the context.


\textbf{[CRITERIA]:} The paper context supports (entails) the citation sentence. In cases where more than one paper is referenced in the sentence, as long as context in which given citation number fits the paper content, it should be count as entailment as well. In multiple citation cases, the paper does not have to entail whole sentence. Scoring rubric is as follows:
0: The given context does not align with the context in which the citation number appears in the sentence, or with the overall meaning of the sentence itself.
1: The given context is compatible with the context in which the citation number appears in the sentence.





[CONTEXT]: In this paper, we learn a temporal, multisensory representation that fuses the visual and audio components of a video signal. We propose to train this model without using any manually labeled data. That is, rather than explicitly telling the model that, e.g., it should associate moving lips with speech or a thud with a bouncing ball, we have it discover these audio-visual associations through self-supervised training [3] . Specifically, we train a neural network on a "pretext" task of detecting misalignment between audio and visual streams in synthetically-shifted videos. The network observes raw audio and video streams -some of which are aligned, and some that have been randomly shifted by a few seconds -and we task it with distinguishing between the two. This turns out to be a challenging training task that forces the network to fuse visual motion with audio information and, in the process, learn a useful audio-visual feature representation.
We demonstrate the usefulness of our multisensory representation in three audiovisual applications: (a) sound source localization, (b) audio-visual action recognition; and (c) on/off-screen sound source separation. Figure 1 shows examples of these applications. In Fig. 1(a) , we visualize the sources of sound in a video using our network's learned attention map, i.e. the impact of an axe, the opening of a mouth, and moving hands of a musician. In Fig. 1(b) , we show an application of our learned features to audio-visual action recognition, i.e. classifying a video of a chef chopping an onion. In Fig. 1(c) , we demonstrate our novel on/off-screen audio source separation model's ability to separate the speakers' voices by visually masking them from the video. The main contributions of this paper are: 1) learning a general video representation that fuses audio and visual information; 2) evaluating the usefulness of this representation qualitatively (by sound source visualization) and quantitatively (on an action recognition task); and 3) proposing a novel video-conditional source separation method that uses our representation to separate on-and off-screen sounds, and is the first method to work successfully on real-world video footage, e.g. television broadcasts. Our feature representation, as well as code and models for all applications are available online.

CITATION NUMBER: 50


\textbf{[ANSWER]:} Video-audio joint analysis [49], [50], [51] utilizes modality consistency to learn representations.
}

% ---------- Remaining tasks ----------
\newcommand{\UserPositioningCheckReal}{%
\textbf{[QUERY]:} Your task is to write a related work section paragraph for a scientific paper that states the paper's contribution or its position among the literature in each paragraph.


\textbf{[CRITERIA]:} The generated paragraph should explicitly or implicitly mention the main paper's contribution or position among existing literature. Ensure that contribution and/or positioning statements should be aligned with the specific focus of the paragraph. Scoring rubric is as follows:
0: The answer paragraph fails to explicitly or implicitly mention the main paper's contribution or position among existing literature.
1: The answer paragraph explicitly or implicitly mention the main paper's contribution or position among existing literature.





\textbf{[ANSWER]:} Averaging of model parameters has a long history in maximum-likelihood estimation. Uniform averaging and exponential moving averaging (EMA) of parameters have been employed in a variety of settings to reduce variance and improve generalization [7--10]. While these techniques are well understood for convex optimization problems, generative adversarial networks correspond to saddle-point (minimax) problems, and the theoretical guarantees for averaging in the convex case do not directly carry over to the adversarial setting.
}

\newcommand{\UserPositioningTypeReal}{%
\textbf{[QUERY]:} Your task is to write related work section for a scientific paper
that states the paper's contribution or its position among the literature in
each paragraph.


\textbf{[CRITERIA]:} In a multiple paragraph related work section, every single
paragraph should state the contribution of the paper and/or its position
among the literature. Ensure that contribution and/or positioning statements
should be aligned with the specific focus of each paragraph. Scoring rubric is
as follows:
0: The answer fails to state contribution of the paper and/or its position among the literature in each related-work paragraph.
1: The answer states contribution of the paper and/or its position among the literature in each related-work paragraph.





\textbf{[ANSWER]:} Marginal-likelihood approximations have long been a cornerstone of scalable Gaussian-process inference, predating the introduction of pseudo-inputs and the variational-inducing-point framework. Early stochastic treatments of the evidence lower bound (e.g., Hernández-Lobato \& Adams, 2015; Salimbeni \& Deisenroth, 2017) can be seen as precursors to the Bayesian perspective adopted in our work. While those methods approximate the ELBO by subsampling data points, they still rely on a fixed set of inducing variables that must be learned jointly with the latent representation. In contrast, the stochastic active-set (SAS) approach we propose replaces learned inducing points with randomly sampled subsets of the training data and derives a stochastic estimator of the log-marginal likelihood that remains unbiased. This positions SAS as a true Bayesian counterpart to earlier stochastic ELBO approximations, offering a simpler optimisation landscape and eliminating the need for explicit inducing-point optimisation.  

Active-set approximations were among the first sparse GP techniques (Smola \& Scholkopf, 2001) and have been empirically shown to remain stable even when the subset size is a small fraction of the dataset (Quiñonero-Candela \& Rasmussen, 2005). However, subsequent work reported that naïve random selection of active sets can induce noisy, non-smooth objective fluctuations that hinder gradient-based optimisation, motivating the shift toward continuous pseudo-inputs (Snelson \& Ghahramani, 2006). Our contribution directly addresses this limitation: by integrating the active-set selection into a stochastic optimisation loop, SAS treats the subset as a random variable and averages over many realizations, thereby smoothing the training objective without sacrificing the computational benefits of small subsets. Empirically, this yields faster convergence and more reliable training of GP decoders than both fixed-active-set schemes and variational inducing-point methods.  

The relationship between cross-validation (CV) and Gaussian-process marginal-likelihoods has been explored for model selection (Wahba, 1990; Rasmussen, 2006), and recent theoretical work has shown an exact equivalence between exhaustive CV and the log-marginal likelihood (Fong \& Holmes, 2020). Existing GP-LVM extensions that incorporate back-constraints or amortised inference (Damianou et al., 2016; Sun et al., 2022) still rely on deterministic inducing variables and therefore inherit the optimisation challenges associated with them. By leveraging the CV--marginal-likelihood equivalence, our SAS framework obtains a tractable stochastic estimator of the marginal likelihood that naturally accommodates random active-set sampling. This connection enables us to train GP decoders that discover richer latent structures, scale to larger datasets, and outperform both variational auto-encoders and state-of-the-art GP-LVMs that use fixed inducing points.
}

\newcommand{\UserActionabilityReal}{%
\textbf{[QUERY]:} Your task is to write a review comment for a scientific paper. The comment should be actionable. Those actions should be clearly identifiable and concrete.


\textbf{[CRITERIA]:} Explicit actions or suggestions are direct or apparent. Authors can directly identify modifications they should apply to their draft. Clarification questions should be treated as explicit statements if they give a direct action. However, implicit actions need to be inferred from the comment. This includes missing parts that need to be added. Authors can deduce what needs to be done after reading the comment. For concrete actions, the authors know exactly what needs to be done and how to apply the action. However, for vague actions the authors still don't know how to carry out this action. Scoring rubric is as follows:
1: The comment lacks meaningful information to help authors improve the paper. Authors do not know what they should do after reading the comment.
2: The comment includes an implicitly stated action or an action that can be inferred. However, the action itself is vague and lacks detail on how to apply it.
3: The comment explicitly states an action but is vague on how to execute it.
4: The comment implicitly states an action but concretely states how to implement the inferred action.
5: The comment contains an explicit action and concrete details on how to implement it. Authors know exactly how to apply it.


\textbf{[ANSWER]:} * The algorithm still requires ground-truth annotations of the datasets (even if it doesn't require a CoT rationale).
}

\newcommand{\UserGroundingReal}{%
\textbf{[QUERY]:} Your task is to write a review comment for a scientific paper. The comment should refer to a specific part of the paper and clearly identify the issue with that part.


\textbf{[CRITERIA]:} For fully grounded comment, the author can accurately pinpoint the section, table, figure, or unique aspect being addressed. For weak grounded comment, the author can make an educated guess but cannot precisely identify the referenced part. For specificity, the comment should detail what is wrong or missing in the referenced part. If external work is mentioned, it should also provide specific examples. Scoring rubric is as follows:
1: The comment is not grounded at all. It does not identify a specific area in the paper. The comment is highly unspecific.
2: The authors cannot confidently determine which part the comment addresses. Further, the comment does not specify what needs to be addressed in this part.
3: The authors cannot confidently determine which part the comment addresses. However, the comment clearly specifies what needs to be addressed in this part.
4: The comment explicitly mentions which part of the paper it addresses, or it should be obvious to the authors. However, this comment does not specify what needs to be addressed in this part.
5: The comment explicitly mentions which part of the paper it addresses, and it is obvious to the authors. The comment specifies what needs to be addressed in this part.


\textbf{[ANSWER]:} - fetched toxic sentences for rewriting (line 215), - toxicity classifier (line 218) - annotator bias These choices introduce biases into the datasets (unavoidable). The paper would benefit from attempting to quantify how useful a model trained on ParaDetox is in a more general setting. e.g. performance on other (non-parallel) toxicity datasets.
}

\newcommand{\UserVerifiabilityReal}{%
\textbf{[QUERY]:} Your task is to write a review comment for a scientific paper. The claims in your comment should be justified or proved by providing logical reasoning, using common sense, or referencing external sources.


\textbf{[CRITERIA]:} Claim justification-verification can be done either by logical reasoning supporting the claim, common sense knowledge in the field verifying the claim (e.g., referencing established practices or standards), or external references substantiating the claim. Scoring rubric is as follows:
1: The comment contains a claim without any supporting evidence or justification.
2: The comment provides some support for its claim, but the justification is vague, insufficient, or not fully articulated. Authors may struggle to follow the reasoning.
3: The comment provides support for its claim, but key elements are missing, such as specific examples, detailed explanations, or supporting references. Authors must make a significant effort to follow the justification.
4: The comment's claim is sufficiently supported but has minor gaps. The reviewer could provide a more detailed explanation or reference.
5: The claim is thoroughly supported by explicit, sufficient, and robust evidence. This can be achieved through: - Clear and precise reasoning or explanation. - Specific and relevant references to external works or data. - Logical and unassailable common-sense arguments.


\textbf{[ANSWER]:} 5. KDSM essentially is matching and retrival. The global matching matrix will cause K-k invalid text prompts, which would waste computation.
}

\newcommand{\UserHelpfulnessReal}{%
\textbf{[QUERY]:} Your task is to write a review comment for a scientific paper. The comment should be useful for the authors to help improving the paper.


\textbf{[CRITERIA]:} A helpful review should be actionable, grounded on a specific part of the paper, provide justification or evidence to its claims. Scoring rubric is as follows:
1: The comment fails to identify meaningful weaknesses or suggest improvements, leaving the authors with no actionable feedback.
2: The comment identifies a weakness or improvement area but is vague, lacks clarity, or provides minimal guidance, making it only slightly beneficial for the authors.
3: The comment identifies weaknesses or areas for improvement but is incomplete or lacks depth. While the authors gain some insights, the feedback does not fully address their needs for improving the draft.
4: The comment provides clear and actionable feedback on weaknesses and areas for improvement, though it could be expanded or refined to be fully comprehensive and impactful.
5: The comment thoroughly identifies weaknesses and offers detailed, actionable, and constructive suggestions that empower the authors to significantly improve their draft.


\textbf{[ANSWER]:} • Is Lemma 1 a well-known result for any DAG? The Jacobian can be transformed to an upper or lower triangular matrix with unit diagonals by permuting the nodes and hence the determinant is trivially 1. I might have missed something here.
}

% ============================================================
% 排版顺序控制（\TaskUserPromptFigure 已不含 \clearpage，
% 此处显式控制分页：RU 两两合并到一页；RW 各占一页，
% Positioning Check 放在最后。）
% ============================================================

% ---------- RU 第 1 页：Actionability + Grounding Specificity ----------
\clearpage
\TaskUserPromptFigure
  {RU -- Actionability}
  {Complete user prompt (test instance \texttt{actionability\_test\_0001}) for Review Utility actionability.}
  {fig:user-ru-actionability}
  {\UserActionabilityReal}

\TaskUserPromptFigure
  {RU -- Grounding Specificity}
  {Complete user prompt (test instance \texttt{grounding\_specificity\_test\_0001}) for Review Utility grounding specificity.}
  {fig:user-ru-grounding}
  {\UserGroundingReal}

% ---------- RU 第 2 页：Verifiability + Helpfulness ----------
\clearpage
\TaskUserPromptFigure
  {RU -- Verifiability}
  {Complete user prompt (test instance \texttt{verifiability\_test\_0001}) for Review Utility verifiability.}
  {fig:user-ru-verifiability}
  {\UserVerifiabilityReal}

\TaskUserPromptFigure
  {RU -- Helpfulness}
  {Complete user prompt (test instance \texttt{helpfulness\_test\_0001}) for Review Utility helpfulness.}
  {fig:user-ru-helpfulness}
  {\UserHelpfulnessReal}

% ---------- RW：Coherence Part 1 ----------
\clearpage
\TaskUserPromptFigure
  {RW -- Coherence (Part 1)}
  {Complete user prompt (test instance \texttt{coherence\_0001}) for Related Work coherence, first half of the context (Part 1).}
  {fig:user-rw-coh-p1}
  {\UserCoherenceRealPartOne}

% ---------- RW：Coherence Part 2 ----------
\clearpage
\TaskUserPromptFigure
  {RW -- Coherence (Part 2)}
  {Complete user prompt (test instance \texttt{coherence\_0001}) for Related Work coherence, second half of the context together with the citation number and answer (Part 2).}
  {fig:user-rw-coh-p2}
  {\UserCoherenceRealPartTwo}

% ---------- RW：Positioning Type ----------
\clearpage
\TaskUserPromptFigure
  {RW -- Positioning Type}
  {Complete user prompt (test instance \texttt{positioning\_type\_0001}) for Related Work positioning type.}
  {fig:user-rw-pos-type}
  {\UserPositioningTypeReal}

% ---------- RW：Positioning Check（放在最后） ----------
\clearpage
\TaskUserPromptFigure
  {RW -- Positioning Check}
  {Complete user prompt (test instance \texttt{positioning\_check\_0001}) for Related Work positioning consistency.}
  {fig:user-rw-pos-check}
  {\UserPositioningCheckReal}

% ============================================================
% 第三方模型 prompt 浮动体：rationale 质量评测（每轮一图）
% ============================================================

\clearpage
\begin{figure*}[t]
\begin{promptboxbig}{Rationale-quality judge -- Round-1}
  \ptext{\JudgeQualityRoundOne}
\end{promptboxbig}
\caption{Complete Round-1 judge prompt. Rationale pairs come from SO and RS
predictions (50 pairs per task, 7 tasks); primary judges are
\texttt{glm-5.3-flash} and \texttt{doubao-seed-2.0-lite}, with
\texttt{MiniMax-M3} as adjudicator.}
\label{fig:judge-r1}
\end{figure*}

\clearpage
\begin{figure*}[t]
\begin{promptboxbig}{Rationale-quality judge -- Round-2}
  \ptext{\JudgeQualityRoundTwo}
\end{promptboxbig}
\caption{Complete Round-2 judge prompt for rationale--score consistency.
The material block is identical to Round-1 (Figure~\ref{fig:judge-r1}).}
\label{fig:judge-r2}
\end{figure*}

% ============================================================
% 第三方模型 prompt 浮动体：DS--RF 过渡审计（两页两图）
% ============================================================

% ---------- 第 1 页：System prompt + 材料块 ----------
\clearpage
\begin{figure*}[t]
\begin{promptboxbig}{DS--RF transition audit (Part 1)}
  \ptext{\TransitionAuditPartOne}
\end{promptboxbig}
\caption{First part of the DS--RF transition audit prompt: the system
message and the evaluation material block. The task instruction and scoring
criteria are shared across the seven in-domain tasks; sample and transition
types identify the DS--RF disagreement category under audit.}
\label{fig:transition-audit-p1}
\end{figure*}

% ---------- 第 2 页：判定说明 + 错误类型 + JSON ----------
\clearpage
\begin{figure*}[t]
\begin{promptboxbig}{DS--RF transition audit (Part 2)}
  \ptext{\TransitionAuditPartTwo}
\end{promptboxbig}
\caption{Second part of the DS--RF transition audit prompt: the score-support
judgment options, the sentence-level error taxonomy for harmful transitions,
and the required JSON output structure. Continues from
Figure~\ref{fig:transition-audit-p1}.}
\label{fig:transition-audit-p2}
\end{figure*}

% ============================================================
% 第三方模型 prompt 浮动体：rationale 修正编辑器（两页两图）
% ============================================================

% ---------- 第 1 页：System prompt + 材料块 ----------
\clearpage
\begin{figure*}[t]
\begin{promptboxbig}{Rationale correction (Part 1)}
  \ptext{\RationaleCorrectionPartOne}
\end{promptboxbig}
\caption{First part of the rationale-correction prompt: the system message
and the material block. The material is a JSON object containing the task,
criteria, evaluated answer, the original rationale, and a list of
sentence-level error targets with correction guidance.}
\label{fig:rationale-correction-p1}
\end{figure*}

% ---------- 第 2 页：修正指令 + JSON 输出结构 ----------
\clearpage
\begin{figure*}[t]
\begin{promptboxbig}{Rationale correction (Part 2)}
  \ptext{\RationaleCorrectionPartTwo}
\end{promptboxbig}
\caption{Second part of the rationale-correction prompt: the rewriting
instructions and the required JSON output structure. Continues from
Figure~\ref{fig:rationale-correction-p1}.}
\label{fig:rationale-correction-p2}
\end{figure*}
